\section*{अध्याय \ref{ch:b2:continuous-reactors} --- सतत रिएक्टर और औद्योगिक प्रक्रम}

\begin{solution}{exo:b2:continuous-reactors:1}
$\tau = V/Q_v = 2000/0.50 = \qty{4.0e3}{s}$, अर्थात् \qty{1.1}{h}।
\end{solution}

\begin{solution}{exo:b2:continuous-reactors:2}
$k\tau = 1.2 \times 10^{-3} \times 1800 = 2.16$। टंकी: $X = 2.16/3.16 = 0.68$;
नली: $X = 1 - \mathrm e^{-2.16} = 0.88$।
\end{solution}

\begin{solution}{exo:b2:continuous-reactors:3}
नली: $\tau = \ln 100/0.020 = \qty{230}{s}$। टंकी: $X/(1 - X) = 99 = k\tau$,
$\tau = \qty{4950}{s}$, बीस गुना से अधिक लंबा: उच्च \omterm{def:b2:continuous-reactors:conversion}{रूपांतरण} पर
विलोडित टंकी निकास सांद्रता पर काम करने की भारी क़ीमत चुकाती है।
\end{solution}

\begin{solution}{exo:b2:continuous-reactors:4}
प्रति लीटर: $120 \times 0.50 = \qty{60}{kJ}$ मुक्त, जो
\qty{4.2}{kJ/K} को गर्म करती है: $\Delta T_{\text{ad}} = \qty{14}{K}$। शीतलन की विफलता
मिश्रण को स्पष्ट रूप से गर्म करेगी, पर ख़तरनाक रूप से नहीं, जब तक विलायक
अपने क्वथनांक के निकट न हो या उस परिसर में कोई दूसरी अभिक्रिया आरंभ न हो।
\end{solution}

\begin{solution}{exo:b2:continuous-reactors:5}
टंकी: $Da = X/(1 - X)^2 = 0.8/0.04 = 20$, $\tau = 20/0.010 = \qty{2000}{s}$।
नली: $Da = X/(1 - X) = 4$, $\tau = \qty{400}{s}$। अनुपात 5, जबकि कोटि 1 के लिए $4/\ln 5 =
2.5$: कोटि जितनी ऊँची, विलोडित टंकी अपनी निम्न सांद्रता से उतनी ही
अधिक प्रभावित होती है।
\end{solution}

\begin{solution}{exo:b2:continuous-reactors:6}
$k\tau = 5.0$। तीन टंकियाँ: $X = 1 - (1 + 5/3)^{-3} = 0.947$। एक टंकी: $5/6 =
0.833$। नली: $1 - \mathrm e^{-5} = 0.993$।
\end{solution}

\begin{solution}{exo:b2:continuous-reactors:7}
$X = 0.70$ का अर्थ $k\tau = 0.70/0.30 = 2.33$ है। प्रवाह दुगुना करने से $\tau$ आधा हो जाता है:
$k\tau = 1.17$, $X = 0.54$। 70\,\% वापस पाने के लिए \omterm{def:b2:continuous-reactors:residence-time}{निवास काल}, और इसलिए
आयतन, दुगुना करना होगा।
\end{solution}

\begin{solution}{exo:b2:continuous-reactors:8}
मुक्त: $80 \times 2.0 \times 1.0 \times 0.90 = \qty{144}{kW}$। \qty{330}{K} से \qty{350}{K} तक
गर्म हुई निकास धारा द्वारा ले जाई गई: $4.0 \times 1.0 \times 20 =
\qty{80}{kW}$। जैकेट को अंतर, \qty{64}{kW}, हटाना होगा।
\end{solution}

\begin{solution}{exo:b2:continuous-reactors:9}
रैखिक विमाएँ $\times 10$: आयतन $\times 1000$, दीवार का क्षेत्रफल $\times 100$, प्रति
इकाई आयतन क्षेत्रफल $\div 10$। मुक्त ऊष्मा आयतन के साथ बढ़ती है, दीवार से होकर
हटाई गई ऊष्मा उसके क्षेत्रफल के साथ: बड़ा रिएक्टर, समान ताप-अंतर पर, प्रति लीटर
दस गुना कम ऊष्मा हटाता है।
\end{solution}

\begin{solution}{exo:b2:continuous-reactors:10}
नली: $\tau = \ln(0.05/0.10)/(0.05 - 0.10) = \qty{13.9}{min}$, $c_B/c_{A,0} =
\frac{0.10}{-0.05}(\mathrm e^{-1.386} - \mathrm e^{-0.693}) = 0.50$। टंकी: $\tau =
1/\sqrt{0.005} = \qty{14.1}{min}$, $c_B/c_{A,0} = 1.414/(2.414 \times 1.707) =
0.34$। टंकी में ताज़ा A और तैयार B एक ही आयतन में मिले रहते हैं: कुछ
B सदा इतनी देर रुकता है कि C में बदल जाए, जबकि कुछ A अभी-अभी
पहुँचा है; नली में हर खंड का इतिहास समान है।
\end{solution}

\begin{solution}{exo:b2:continuous-reactors:11}
शीतलक (और फ़ीड) का ताप बढ़ाने से निष्कासन रेखा दाईं ओर
खिसकती है। ठंडा प्रतिच्छेद धीरे-धीरे ऊपर उठता है जब तक रेखा S के निचले मोड़ की
स्पर्शरेखा न बन जाए; उसके आगे ठंडी अवस्था लुप्त हो जाती है और टंकी कूदकर
गर्म शाखा पर पहुँचती है (प्रज्वलन)। ताप फिर से घटाने पर टंकी तब तक
गर्म शाखा पर रहती है जब तक रेखा ऊपरी मोड़ की स्पर्शरेखा न बन जाए, और केवल
तभी वापस गिरती है (शमन), प्रज्वलन से कम ताप पर: एक
शैथिल्य पाश।
\end{solution}

\begin{solution}{exo:b2:continuous-reactors:12}
बैच: \qty{2.0}{mol/L} का 75\,\% शेष, $1.5 \times 100 = \qty{150}{kJ/L}$,
कार्य-ताप से $\Delta T = 150/4.0 = \qty{37.5}{K}$ ऊपर, जो किसी क्वथनांक
या विघटन तक पहुँचने के लिए पर्याप्त है। अर्ध-बैच: अधिकतम
\qty{0.10}{mol/L} बिना अभिक्रिया के, $\Delta T = 10/4.0 = \qty{2.5}{K}$। धीमा योग
रिएक्टर में संचित ऊर्जा की सीमा बाँध देता है।
\end{solution}

\begin{solution}{pb:b2:continuous-reactors:1}
\textbf{1.} $r = k[\text{एस्टर}][\ce{OH-}]$; समान फ़ीड और 1 : 1 की
रससमीकरणमिति दोनों सांद्रताओं को समान रखती हैं, $r = kc^2$।
\textbf{2.} $1/c - 1/c_0 = kt$, जहाँ $c = 0.10c_0$: $kc_0t = 9$, $t = 9/(0.11
\times 0.10) = \qty{818}{s}$, लगभग \qty{14}{min}।
\textbf{3.} $t_{1/2} = 1/(kc_0) = \qty{91}{s}$; कोटि 2 के लिए दर
सांद्रता के वर्ग के साथ गिरती है, इसलिए उसे आधा करने का समय इस पर निर्भर करता है कि
आरंभ कहाँ से हुआ।
\textbf{4.} $Da = kc_0\tau = 0.011\,\tau$ ($\tau$ सेकंड में)।
\textbf{5.} $Q_vc_0 - Q_vc - kc^2V = 0$, अर्थात् $c_0 - c = k\tau c^2$।
\textbf{6.} $c = c_0(1 - X)$ के साथ: $c_0X = k\tau c_0^2(1 - X)^2$।
\textbf{7.} $Da = 0.90/0.10^2 = 90$।
\textbf{8.} $\tau = 90/0.011 = \qty{8.2e3}{s}$, \qty{2.3}{h}।
\textbf{9.} $V = 0.50 \times 8182 = \qty{4.1e3}{L}$, \qty{4.1}{m^3}।
\textbf{10.} निकास सांद्रता, $0.010\,\unit{mol/L}$: दर
हर जगह प्रवेश की तुलना में सौ गुना कम है।
\textbf{11.} $Da = X/(1 - X) = 9$, $\tau = \qty{818}{s}$, $V = \qty{409}{L}$।
\textbf{12.} नली के लिए दस गुना कम आयतन।
\textbf{13.} $c_{j-1} - c_j = k\tau_1c_j^2$, जहाँ $\tau_1$ एक टंकी का
\omterm{def:b2:continuous-reactors:residence-time}{निवास काल} है।
\textbf{14.} $\tau_1 = 850/(3 \times 0.50) = \qty{567}{s}$, $k\tau_1 =
\qty{62.4}{L/mol}$। $62.4c_j^2 + c_j - c_{j-1} = 0$ हल करने पर: $c_1 = 0.0328$,
$c_2 = 0.0163$, $c_3 = \qty{0.0100}{mol/L}$: 90\,\% \omterm{def:b2:continuous-reactors:conversion}{रूपांतरण}।
\textbf{15.} नली, या श्रेणीबद्ध टंकियाँ (एक टंकी से दस और पाँच गुना
छोटी)। फिर भी एक अकेली टंकी अपने एकसमान ताप, आसान सफ़ाई और
नियंत्रण के लिए, और इसलिए कि यहाँ आयतन सस्ता है, जीत सकती है।
\textbf{16.} $\Delta T_{\text{ad}} = 55\,000 \times 100/(997 \times 4180) =
\qty{1.3}{K}$ (90\,\% पर \qty{1.2}{K})।
\textbf{17.} $55 \times 0.50 \times 0.10 \times 0.90 = \qty{2.5}{kW}$।
\textbf{18.} नहीं: निकास धारा लगभग एक केल्विन की वृद्धि के साथ ऊष्मा
ले जाती है; कोई अनियंत्रण संभव नहीं।
\textbf{19.} बीस गुना अधिक सांद्र: रुद्धोष्म वृद्धि
\qty{26}{K} और ऊष्मा \qty{50}{kW} हो जाती है, जो शीतलन कुंडली के योग्य है; दर $kc_0$
बीस गुना ऊँची होने से 90\,\% के लिए टंकी बीस गुना छोटी होगी
(\qty{0.20}{m^3})।
\textbf{20.} $0.50 \times 0.10 \times 0.90 = \qty{0.045}{mol/s}$, $\times 86\,400
\times 82.03 = \qty{319}{kg}$ सोडियम एथेनोएट प्रति दिन।
\textbf{21.} दर स्थिरांक ताप के साथ बढ़ता है (आरेनियस), इसलिए
आयतन घटता है; क़ीमत है फ़ीड का तापन और, दूसरी
अभिक्रियाओं के लिए, कम \omterm{def:b2:continuous-reactors:selectivity}{वरणात्मकता} और अधिक वाष्प।
\textbf{22.} अकेली विलोडित टंकी का आयतन $\boldsymbol{V \approx
\qty{4.1}{m^3}}$ होना चाहिए।
\end{solution}
